% Generated draft: requires mathematical and source review.
\documentclass[11pt]{article}
\usepackage[margin=25mm]{geometry}
\usepackage{fontspec}
\setmainfont{Latin Modern Roman}
\usepackage{amsmath,amssymb,mathtools,unicode-math}
\setmathfont{Latin Modern Math}
\usepackage{enumitem,needspace,xurl,hyperref}
\setlength{\parindent}{0pt}
\setlength{\parskip}{5pt}
\setlength{\emergencystretch}{3em}
\newcommand{\N}{\mathbb{N}}
\newcommand{\Z}{\mathbb{Z}}
\newcommand{\Q}{\mathbb{Q}}
\newcommand{\R}{\mathbb{R}}
\newcommand{\C}{\mathbb{C}}
\newcommand{\F}{\mathbb{F}}
\newcommand{\A}{\mathbb{A}}
\newcommand{\PP}{\mathbb{P}}
\newcommand{\E}{\mathbb{E}}
\newcommand{\eps}{\varepsilon}
\newcommand{\dd}{\,\mathrm{d}}
\begin{document}
\section*{UnsolvedMath ID: 30005212}
\textbf{Source problem code:} catalogue entry
\subsection{Definitions and mathematical statement}
Let \(n\in\N\) with \(n\geq 3\). The supplied record names nonelliptic differential-inclusion sets \(\Gamma_n\) and asks whether the condition
\[
Dw(x)\in\Gamma_n
\]
when imposed ``sufficiently locally,'' forces \(Dw\) to be ``rigid.'' Here \(Dw\) denotes the derivative of a map \(w\), but the record does not specify the domain or codomain of \(w\), its regularity, the sense of differentiation, or the points at which the inclusion must hold. It also supplies no definition of \(\Gamma_n\), no meaning of ``sufficiently locally,'' and no definition of rigidity. Thus the intended question concerns every integer \(n\geq3\), but a precise quantified implication cannot be stated from the supplied evidence. In particular, rigidity must not silently be identified with constancy of \(Dw\) or affinity of \(w\).

The original extract begins with the continuation question, ``More generally is it true for \(\Gamma_n\) where \(n\in\N\) and \(n\geq3\)?'' Its antecedent is absent. The extract then starts a separately headed problem, ``Problem 3. (Threshold regularity of the eikonal equation).'' That neighboring material must not be substituted for the missing higher-index inclusion formulation.

For context, the neighboring fragment considers
\[
m:\Omega\longrightarrow S^1,
\qquad S^1=\{v\in\R^2:|v|=1\},
\qquad \operatorname{div}m=0\quad\text{in }\mathcal D'(\Omega),
\]
and an entropy condition
\[
\operatorname{div}(\Phi(m))\in L^p_{\mathrm{loc}}(\Omega)
\qquad\text{for every }\Phi\in\mathrm{ENT}.
\]
The printed entropy definition can be read formally as requiring \(\Phi:S^1\to\R^2\) and
\[
\frac{\dd}{\dd t}\Phi(e^{it})\cdot e^{it}=0,
\]
with \(e^{it}\) identified with \((\cos t,\sin t)\). The function class needed to interpret this derivative is not specified in the supplied extract. Its conjectural conclusion is that \(m\) is rigid, again without a supplied definition. The conjecture's range of \(p\), domain hypotheses, and relationship to the \(\Gamma_n\) question remain unresolved. These contextual formulas do not complete the main problem.
\subsection{Short English statement}
For every integer \(n\geq3\), the record asks whether sufficiently local membership of a derivative in \(\Gamma_n\) forces rigidity. The sets, admissible maps, locality condition, and rigidity conclusion are undefined in the supplied evidence, so this remains an incomplete statement requiring mathematical review. The adjacent eikonal conjecture is a separate source fragment.
\subsection{Sources}
\begin{itemize}
\item Supplied local UnsolvedMath registry record, permanent ID \(30005212\), source problem code \texttt{catalogue entry}, in the Oberwolfach Reports Open Problems collection. The project identifies \texttt{data/unsolvedmath/problems.json} and the supplementary registry \texttt{catalog/problems.json}; their contents were not independently inspected for this draft.
\item Supplied source citation: \emph{Calculus of Variations (2023)}, Oberwolfach Reports, DOI 10.4171/owr/2022/37, pp.~2129--2193. \url{https://doi.org/10.4171/owr/2022/37}. The record separately gives proposed year 2022; this bibliographic year discrepancy is unresolved. No individual proposer is supplied.
\item Supplied literature source: Lamy--Lorent--Peng, ``non-elliptic curves.'' \url{https://arxiv.org/abs/2404.02121}. The record does not supply a complete bibliographic citation.
\item Supplied literature source: ``original Oberwolfach report.'' \url{https://doi.org/10.4171/owr/2022/37}.
\item Literature assessment dated exactly \texttt{2026-08-22}: the supplied curation labels the problem open and states that available exact-inclusion regularity results concern different hypotheses and do not establish the asserted higher-index statement for all \(n\). This is supplied dated curation, not an independently verified literature assessment. The record describes it as literature triage derived from OWR curation and requests independent verification of time-sensitive status and cited claims.
\end{itemize}
\end{document}
